\begin{annexitem}{annex:p12.ac.multisets}{thm:p12-ac-multisets}{Nonempty free commutative semigroups and exact cancellation}
\begin{proof}
Flatten f recursively and count each generator occurrence. Associativity
changes only bracketing and commutativity only order. Equal nonempty
multisets can be reordered and rebracketed to identical terms, so
they are AC-equal. Conversely the nonzero finite-support vectors over
the generator alphabet, with vector addition as f, are a model of AC.
Interpreting generators as unit vectors proves that distinct counts
cannot be AC-equal. A positive vector is represented by its sole
atom or a binary fold; the zero vector has no term representative.

Under a simultaneous substitution, each occurrence contributes the
count vector of its one image. Summing those contributions gives the
matrix multiplication law. Counts never cancel against negative ones.
If an identical flattened argument occurs on both sides, its substituted
count vector is the same on both sides. Subtracting that shared vector
from a vector equality preserves exactly the equality. This also proves
the reverse implication by adding it back. Cancellation may leave an
empty balance on a side; that is a numerical comparison of zero
vectors, not an AC-term denoting zero. Original variable images still
have positive total count. Repeating this finite cancellation proves
the stated solution-preservation assertion.
\end{proof}
\end{annexitem}

\begin{annexitem}{annex:p12.ac.single-bound}{thm:p12-ac-single-bound}{The source's positive-coefficient bound and finite minimal search}
\begin{proof}
Both coefficient lists are finite and positive, so K and every
coordinate search bound are finite. Induct on the balance k of a
nonnegative solution. Zero is the empty sum. A solution with k at
most K is itself one of the enumerated solutions. For k greater
than K, some \(a_ix_i>L\): otherwise the left sum would be at
most mL, hence at most K. Likewise some \(b_jy_j>L\).
Put \(q=\operatorname{lcm}(a_i,b_j)\leq L\). The solution
having just \(x_i=q/a_i,y_j=q/b_j\) positive is componentwise
below the original solution. Subtract it. The residual is a
nonnegative solution of balance k-q, strictly smaller than k;
the subtracted solution itself has balance at most K. The induction
therefore generates every solution by solutions inside the finite bound.

If a componentwise minimal nonzero solution had balance above K,
the just constructed subtracted solution and its nonzero residual
would contradict minimality. Thus every global minimum is in the
search. Conversely a searched solution that is not globally minimal
has a smaller nonzero solution, whose balance is smaller and still
within K, so the search detects it. Retained minima are therefore
exactly the global minima. For any nonzero solution choose a minimum
below it in its finite componentwise box, subtract, and induct on
the sum of coordinates. This proves generation by the retained set,
not merely finiteness of a trial search.

If a side has no remaining coefficient, positivity forces every
coordinate on the other side to zero. Cancelled variables are
unconstrained original variables, not such forced zero coordinates.
The procedure must still enforce their nonempty term images rather
than declare an internal zero balance a unit term.
\end{proof}
\end{annexitem}

\begin{annexitem}{annex:p12.ac.system-basis}{thm:p12-ac-system-basis}{Cone decomposition and a proved finite integer remainder bound}
\begin{proof}
There are finitely many supports. Rational Gaussian elimination and
clearing denominators compute each one-dimensional positive kernel
ray; divide by the gcd to get its primitive integer vector. Thus
the ray list and B are computable finite data. Prove first that
these rays generate the real nonnegative cone \(Av=0\).
Rational elimination gives the same rank over the rationals and reals,
so the computed kernel dimensions also apply to this real cone.

For nonzero v, let S be its positive support. If the restricted
kernel has dimension one, v is a positive real multiple of the
enumerated primitive ray on S. Otherwise choose an independent
kernel vector w and put
\(z=w-(\sum w_i/\sum v_i)v\). Then z is a nonzero kernel
vector with coordinate sum zero, so it has both signs. Let
\(t_+=\min_{z_i<0}v_i/(-z_i)\) and
\(t_-=\min_{z_i>0}v_i/z_i\), with indices in S. Both are
positive. The kernel vectors \(v^+=v+t_+z\) and
\(v^-=v-t_-z\) are nonnegative, nonzero (their sums equal
\(\sum v_i>0\)), and each has a strictly smaller support.
Furthermore
\[
 v=\frac{t_-}{t_++t_-}v^+
       +\frac{t_+}{t_++t_-}v^-.
\]
Induction on support size decomposes both vectors into the enumerated
rays, proving finite real cone generation. If no rays exist, there
is no nonzero cone vector.

Now take an integer solution v and a real decomposition
\(v=\sum_j\alpha_jr_j\), \(\alpha_j\geq0\).
Subtract \(\sum_j\lfloor\alpha_j\rfloor r_j\).
The remainder b is an integer nonnegative solution and
\(0\leq b_i\leq\sum_j(r_j)_i=B_i\); indeed the positive
fractional contributions are each strictly below their full ray
contribution. Hence every integer solution is an integer sum of
the ray vectors and one of the finitely enumerated box solutions.
The ray vectors themselves lie in the box. This supplies a finite
generating set G with no unbounded oracle search.

Every minimal nonzero integer solution belongs to G: its expression
as a sum of G-elements cannot have two nonzero summands, because
either would be smaller componentwise. A G-element that is not
globally minimal has a smaller nonzero integer solution; decompose
that solution using G, and one nonzero G-summand is smaller than
the original element. Thus the finite comparison pass removes
exactly the nonminimal elements. Retained H is precisely the
global minimal set. Induction by subtracting a minimum in the
finite box beneath any solution proves that H generates all
integer solutions. Empty cone and zero-column cases follow from
the same argument, not from an assumed nonempty basis.
\end{proof}
\end{annexitem}

\begin{annexitem}{annex:p12.ac.particulars}{thm:p12-ac-particulars}{Last-coordinate-one generators are exactly minimal affine solutions}
\begin{proof}
The augmented system is \(Dw-d_ct=0\). Any vector \((w,1)\)
is a nonnegative homogeneous solution and therefore an integer sum
of its minimal generators. Since all last coordinates are nonnegative
integers and sum to one, exactly one used generator has last coordinate
one, with multiplicity one; all other used generators have last
coordinate zero. The latter project to homogeneous D-solutions.
Thus \(w=p+v\), with p from the selected projection and v
homogeneous. If there is no such generator there is no affine solution.

If p were not affine-minimal, a smaller \((p',1)\) would be
a nonzero augmented solution below its supposed minimal generator.
Conversely an affine-minimal p makes \((p,1)\) minimal: a
smaller augmented solution with last coordinate one contradicts
affine minimality, and one with last coordinate zero is a nonzero
homogeneous h below p, so p-h is a smaller affine solution.
These are the only possible last coordinates below one. Hence
the computed projection is exactly the affine minimal set.
When \(d_c=0\), zero is the unique affine minimum and the
augmented unit last-coordinate vector is a generator. This proves
all stated cases through the actual finite homogeneous algorithm.
\end{proof}
\end{annexitem}

\begin{annexitem}{annex:p12.ac.coverage}{thm:p12-ac-coverage}{Every factor image is nonempty, and every target has an admissible branch}
\begin{proof}
Every selected fresh-parameter column is homogeneous and every fixed
constant column solves its affine system. The coordinate lemma
therefore proves soundness. The row test gives a nonempty term for
every original variable, so no numerical zero is smuggled into the
term model. Finite bases, finite particular choices and the finite
power set of H make every branch enumeration and extraction finite.

Fix any target unifier \(\theta\). Its finitely many variable
images contain a finite set of generator atoms \(Q\). For each input
constant \(c\) its image-count column \(w_c\) solves \(Dw_c=d_c\).
Write \(w_c=p_c+\sum_j h_ja_{j,c}\), with the actual minimal
particular \(p_c\). For every other atom \(q\) in \(Q\), its column is
homogeneous, so write \(w_q=\sum_jh_ja_{j,q}\). This includes
an allowed constant absent from the input, whose input offset is zero.
All coefficients are nonnegative integers by the proved generation.
These are constructive representations: choose a finite particular
column below the target, and repeatedly subtract a known minimal
homogeneous column below each nonzero residual. Such a column exists
by generation, and each subtraction strictly decreases its integer
coordinate sum. The recorded multiplicities give the stated coefficients.
Select exactly those j for which some \(a_{j,q}\), including
the constant residual coefficients, is positive. Assign \(z_j\)
the term whose q-counts are those \(a_{j,q}\). A selected j
has positive total count by definition, so this is a genuine
nonempty AC term, not an empty image or a zero-multiplicity variable.
No assignment is needed for an omitted j, which occurs nowhere
in this candidate answer.

For any row \(i\), the target's total image count is positive. If its
particular constant sum is zero, some selected \(h_j\) has positive
row entry, since its coefficient contributes to that target count.
Thus the selected branch passes the row test. Applying its one
finite simultaneous factor gives exactly every target atom count:
the fixed constants contribute \(p_c\) and its fresh images contribute
the generated residuals and all homogeneous atom columns. Hence
\(x_i\theta=_{AC}x_i\sigma_{p,I}\lambda\) on X. Fresh
names may occur in the target itself; their inserted images are
not reapplied during this simultaneous composition. No acyclicity
or uniqueness of the factor is presumed.

For n=0 there is no variable image to construct. Fixed ground
multisets either agree in all equations or do not; identity is
a solution exactly in the former case. On the empty comparison
domain it covers every substitution. This also covers empty input.
Duplicate rows impose the same equations and change neither the
solution columns nor this extraction/coverage argument.
\end{proof}
\end{annexitem}

\begin{annexitem}{annex:p12.ac.minimality}{thm:p12-ac-minimality}{Indecomposable columns force the same selected family}
\begin{proof}
Take two candidates \(\sigma_{p,I}\) and \(\sigma_{q,J}\),
and suppose the first covers the second on X. Each selected source
parameter has a nonempty factor image. Any atom outside the target's
constants and parameters would appear in a row where that nonzero
source h-column is positive, contradicting the target zero count
at that atom. Thus factor images use only those target atoms.

For a target parameter \(j\), its \(h_j\) column equals
\(\sum_{i\in I}h_i a_{i,j}\). A minimal nonzero homogeneous
solution is additively indecomposable: two nonzero summands would
give a smaller nonzero homogeneous solution. Consequently exactly
one summand equals \(h_j\), with multiplicity one. Distinct elements
of \(H\) are distinct vectors, so it must be the source column \(i=j\).
This proves \(J\subseteq I\) and fixes all parameter-column
contributions.

For each constant \(c\), comparison gives
\(q_c=p_c+\sum_{i\in I}h_i b_{i,c}\), hence \(q_c\geq p_c\).
Both are affine minima in the same \(P_c\), so they must be equal.
A nonzero \(h_i\) with a positive coefficient would give a positive
coordinate in that difference; therefore every b coefficient is
zero. A source parameter not in J would now have no constant or
parameter contribution at all, forcing an empty factor image.
That is impossible in AC. Thus \(I\subseteq J\) too, and
\((p,I)=(q,J)\). Distinct candidates are incomparable.

Canonical generation removes only literal vector/name duplicates;
the same proof also excludes hidden instantiation-equivalent
duplicates up to fresh renaming. Completeness has already been
proved, so this is a minimal complete set. If nonminimal numerical
generators were used, the indecomposability step would fail and
the conclusion could not be inferred merely from a finite output.
\end{proof}
\end{annexitem}
